\scanpage{266}
\mspage{7}

formée des $j\in I$ qui ne sont pas parmi les $s_j$ déjà construits, sauf \deleted{\unsure{du}} le dernier $s_i$. \deleted{\unsure{Cet $i$}, \unsure{$I'$}} \inserted{(donc $s_i\in I'$, $\operatorname{card}I'\geq2$)} Alors \deleted{l'une des catégories} $I'$ est isom. à un intervalle \unsure{standard|standart}, \unsure{son|avec} $s_i$ correspondant à l'origine. Une des catégories $i\backslash I$, $I_{/i}$ est ponctuelle, l'autre \deleted{\unsure{équival}} isom. à une catégorie $\mathbb{\Delta}_j$, et de plus,

$i$ est objet initial ou final de $\mathbb{\Delta}_j$. Si c'est l'objet initial $0$, on prend pour $s_{i+1}$ le successeur \inserted{immédiat} de $i$ dans $\mathbb{\Delta}_j$. Si c'est l'objet final $j$, on prend pour $s_{i+1}$ le prédécesseur immédiat $\unsure{j'|j}-1$.

Si on était parti du \unsure{but|bout} \inserted{$\unsure{b|h}$} de $I$, au lieu de partir de la source, on aurait trouvé la relation d'ordre totale opposée. \deleted{Pour cela} \unsure{Affectant|Affectons} l'ancien but du numéro d'ordre $0$, on trouve une renumérotation de $I$, $i$ devenant $n-i$, et les \deleted{\unsure{deux}} flèches de transition qui étaient \deleted{d'un signe} directes pour la \deleted{première} relation d'ordre initiale, deviennent rétrogrades dans la relation \deleted{\unsure{oppos}} \inserted{d'ordre} opposée, et inversement. On voit ainsi :

\uline{Proposition 2} Si $I$ est un intervalle \inserted{standard}, il existe un unique \deleted{intervalle $I'$,} couple $(I',\Theta)$, où $I'$ est un intervalle standard, et $\Theta$ un iso $I\simeq I'$ qui transforme le but de $I$ dans la source de $I'$ (et la source de $I$ en le but de $I'$.

Cet intervalle $I'$ est appelé l'intervalle \uline{opposé} de $I$. Il faut le voir, intrinsèquement, comme la ``même'' catégorie $I$, mais avec une ``orientation'' opposée.

% archive sheet number: 265
